\begin{textsample}{POS Dim 4 – pi – Score 11.00 – pi/pi\_em\_58.txt}  \label{ex:f4_pos_124}
A Área de Ciências Humanas e Sociais Aplicadas e a Integração entre os Componentes O mundo contemporâneo aponta para a construção de uma sociedade intercultural, com vistas à garantia da cidadania e da valorização das diversidades, legitimada pelos princípios da equidade, justiça, igualdade, solidariedade, ética, alteridade e a autonomia dos sujeitos envolvidos no processo de ensino e aprendizagem.
Nessa perspectiva, emerge a necessidade de interlocução das diversas áreas, com foco na concepção de educação integral, que considere os sujeitos em sua integralidade, que vise promover o desenvolvimento dos estudantes nas suas dimensões : intelectual, socioemocional, física e cultural.
Articulada aos propósitos da Constituição Federal ( 1988 ) e da BNCC ( 2018 ), a área de Ciências Humanas e Sociais Aplicadas, cujo foco é a dimensão humanística, visa desenvolver no jovem e no adulto, enquanto sujeitos do Ensino Médio, a sua dimensão humana integral e emancipatória, por meio de aquisição de capacidades voltadas a fortalecer o sentimento de solidariedade, respeito às diferenças, reconhecimento e valorização das diversidades.
Essa finalidade requer um ensino voltado à construção intencional de processos educativos que promovam aprendizagens sintonizadas com as necessidades, com as possibilidades e os interesses dos estudantes.
Nesse aspecto, a aprendizagem deve ser contextualizada com a realidade social local, regional, nacional e com os desafios da sociedade contemporânea, por meio de práticas pedagógicas que favoreça o reconhecimento e valorização das \textbf{identidades} culturais, étnicas e linguísticas.
Tudo isso com vista a uma aprendizagem significativa para a vida, capaz de promover o protagonismo juvenil e a construção do projeto de vida dos sujeitos do Ensino Médio.
A BNCC ( 2018 ) traz algumas inovações importantes na área de ciências humanas e sociais, podendo ser citadas a contextualização do próprio conceito da área, o que marca uma inovação de ordem epistemológica.
Outro aspecto inovador também está relacionado ao tratamento dos temas contextualizados aos cenários locais dos estudantes, com vista a estimular a noção de pertencimento, sua ação protagonista, bem como sua atuação autoral no conhecimento e nos projetos de pesquisa e de intervenção no \textbf{território}.
Para a construção de uma sociedade mais cidadã e de sujeitos conscientes acerca das questões sociais, políticas, culturais, \textbf{religiosa} e econômica é preciso ampliar o “ focus ” e consolidar o papel das Ciências Humanas e Sociais Aplicadas, garantindo o legítimo e necessário espaço de articulação entre todas as ciências, oportunizando a formação de sujeitos comprometidos com a coletividade, com uma nova ordem social que desloque a centralidade do \textbf{capital} para a centralidade ética da vida humana.
Com o reconhecimento da centralidade do \textbf{homem} e da vida humana, as Ciências Humanas e Sociais aplicadas, integradas por Filosofia, Geografia, História e Sociologia se destacam nesse cenário, pois em conformidade com a BNCC ( 2018 ), estas possuem o papel de ampliar e aprofundar “ as aprendizagens essenciais desenvolvidas até o Ensino Fundamental, e consolidá -las no Ensino Médio, sempre orientadas para uma educação ética ” ( BNCC, 2018, p. 547 ).
Nesse contexto, o grande objetivo da área é direcionar os/as estudantes a refletirem sobre a sociedade e seu papel social, cultural, histórico, ético e político sobre o mundo local e global.
A importância desta área é sintetizada na BNCC ( 2018 ) do seguinte modo : para se viver em sociedade, e em cujas bases destacam -se as ideias de justiça, solidariedade e igualdade ” se faz necessário uma proposta que tenha como fundamento a compreensão e o reconhecimento das diferenças, o respeito aos direitos humanos e à interculturalidade, e o combate aos preconceitos ( BNCC, 2018, p. 547 ).
As Ciências Humanas são constituídas pelas categorias : Tempo e Espaço ; Território e \textbf{Fronteira} ; Indivíduo, Natureza, Sociedade, Cultura e Ética ; e Política e Trabalho.
Assim, o estudo dessas categorias é essencial para compreensão da condição humana, e da identificação do indivíduo como ser social, pois, agrega temáticas que transversalizam de modo a alcançar a interdisciplinaridade.
A categoria tempo, central para o estudo da História e das Ciências Humanas e Sociais, permite compreender todas as demais áreas, ao possibilitar análise dos processos marcados por continuidade, mudanças e rupturas.
Assim, tempo deve ser compreendido de acordo com Le Goff ( 1984 ), numa dimensão múltipla que ultrapasse a cronológica, e englobe a dimensão simbólica e abstrata.
Conforme assevera Delgado ( 2020 ) a complexidade integrante à noção de tempo refere -se às temporalidades múltiplas que se enlaçam, uma vez que as experiências vividas e a História em transformação são conformadas por processos e acontecimentos.
Segundo expressa a BNCC, a compreensão de espaço deve ultrapassar suas representações cartográficas e ser contemplado nas dimensões histórica e cultural.
As noções de \textbf{território} e \textbf{fronteiras} são essenciais para compreender os processos identitários dos grupos nas suas organizações culturais, econômicas, \textbf{religiosas}, sociais e políticas.
Nesse processo, a abordagem filosófica é essencial, porque é preciso pensar o cotidiano, apoiados nas diferentes formas do pensamento filosófico integrada às abordagens sociológicas, \textbf{antropológicas}, políticas e econômicas da vida cotidiana.
Essa integração tem o propósito de preparar os educandos para refletir e, assim, interpretar os processos sociais, levando -os a desnaturalizar seu cotidiano.
Dessa forma, será possível conduzir a uma percepção da interculturalidade espaço-temporal de experiências, normas de convívio, condutas e papéis sociais necessários para a vida em sociedade, de forma a contribuir para formar sujeitos mais cidadãos.
O foco da educação integral, além de ensinar conhecimentos intrínsecos de cada Área, habilita o estudante a resolver os problemas da vida e dos desafios contemporâneos, numa articulação entre o conhecimento curricular e os acontecimentos existentes da sociedade ao seu redor.
Nesse aspecto, os componentes da Área de Ciências Humanas e Sociais Aplicadas, cujo foco é formar o discente para exercer sua cidadania ativa, protagonizando seus desejos, necessidades e objetivos de vida, devem considerar as diferentes dimensões que constituem o estudante enquanto indivíduo.
Dessa forma, o trabalho interdisciplinar dos componentes possibilita compreender a formação cognitiva articulada as dimensões socioemocionais, físicas, ambientais, culturais, sociais, éticas e históricas, com vistas a uma formação integral, agregando elementos relacionados as singularidades dos sujeitos e da vida social, considerando as \textbf{identidades} individuais e coletivas.
Nessa perspectiva, visando desenvolver um processo de ensino para a aquisição de competências voltadas ao estímulo da curiosidade, do senso crítico, da autodeterminação e da autonomia dos estudantes, cada componente deverá desenvolver seus objetos do conhecimento de forma articulada, focando as especificidades que norteiam cada componente, seus conceitos e práticas, visto que esses objetos estão em conexão com as competências específicas e habilidades, norteadas por um tema geral.
Assim, o trabalho deve ser permeado pela articulação dos diversos saberes, onde os limites entre cada componente acabam se mesclando.
Dessa forma, torna -se fundamental um planejamento por área, que considere também as 7 ( sete ) categorias definidas pela BNCC que permeiam todos os componentes curriculares da área de Ciências Humanas e Sociais Aplicadas, a saber : Tempo e Espaço ; \textbf{Territórios} e Fronteiras ; Indivíduo, Natureza ; Sociedade ; Cultura e Ética ; Política e Trabalho.
Neste sentido, no processo de definição ou seleção das competências, habilidades, objetivos e sua relação com os objetos do conhecimento, as referidas categorias possibilitam um exercício de escolha de forma mais objetiva, ao considerar a relação entre as competências e habilidades com determinada categoria.
Neste propósito, é necessário o alinhamento dos objetos do conhecimento, seguindo a sequência das competências específicas e suas habilidades, que devem ser trabalhados de forma interdisciplinar e transdisciplinar, contextualizados com a realidade dos discentes e articuladas às bases teóricas, políticas, éticas, ontológicas, epistemológicas e conceituais de cada componente.
Para auxiliar a inserção dos estudantes nas múltiplas dimensões que \textbf{atravessam} a formação integral, a BNCC ( 2018 ) elenca os Temas Contemporâneos Transversais – TCTs, que devem ser desdobrados em todas as áreas do conhecimento e componentes curriculares.
As propostas estão relacionadas à perspectiva do conhecimento holístico e operante e buscam articular os conhecimentos escolares, organizando as ações de ensino, de forma flexível, podendo ser ou não desenvolvidas, dependendo das referências disciplinares preestabelecidas.
Os TCTs são desenvolvidos de modo contextualizado e transversal, e para atender as diferentes demandas, suas abordagens foram divididas em três níveis ascendentes de complexidade, de forma a tratar os TCTs de maneira intradisciplinar, interdisciplinar e transdisciplinar.
Sobre a intradisciplinaridade, está relacionada ao nível de integração dos diferentes objetos do conhecimento abordados em um só componente curricular.
A interdisciplinaridade trata de reunir as possibilidades de produção de conhecimentos que trazem cada componente curricular.
É o trabalho coletivo entre diversas disciplinas com o objetivo de desenvolver uma variedade de conteúdos ao tratar do mesmo assunto ; e por fim, a transdisciplinaridade é a interação geral que permeia em todas as áreas em torno de um mesmo conteúdo.
Sendo assim, as propostas podem ser trabalhadas tanto em um ou mais componentes de forma intradisciplinar, interdisciplinar ou transdisciplinar, mas sempre transversalmente às áreas do conhecimento.
Os temas contemporâneos transversais surgem contextualizados nos objetos do conhecimento da Formação Geral Básica e dos Itinerários Formativos.
São eles : Meio Ambiente, Saúde, Multiculturalismo, Economia, Cidadania e Civismo, Ciência e Tecnologia.
A seguir lista -se alguns direcionamentos de como os componentes de Filosofia, Geografia, História e Sociologia podem explorá -los em suas abordagens.
O tema Meio Ambiente, está subdividido em temas relacionados, a Educação Ambiental e Educação para o Consumo.
Nesta temática deve -se ressaltar as relações \textbf{Homem} x Natureza, numa perspectiva histórica, social, geográfica e filosófica que contribuam para o despertar da consciência ambiental e para o desenvolvimento sustentável, gerando atitudes de preservação e conservação.
Trabalhar as temáticas que envolvem o meio ambiente asseguram que as próximas gerações internalizem a importância de cuidar diariamente do planeta, repensando suas atitudes de consumo, com vistas a reduzir os danos das nossas rotinas produção e descarte.
A temática sobre Saúde tem seus conhecimentos distribuídos em subtemas de Saúde, Educação Alimentar e Nutricional.
Prepara o estudante para o autocuidado, possibilitando a relação entre saúde e meio ambiente.
É uma temática fundamental para trabalhar questões de promoção da saúde e prevenção de doenças, estimulando que os estudantes revejam seus hábitos e atitudes para que possam obter uma maior qualidade de vida, conscientizando e ajudando a conscientizar a sociedade sobre a necessidade de uma vida mais saudável.
O Multiculturalismo aborda temas relacionados a Diversidade Cultural e Educação para valorização do multiculturalismo nas matrizes históricas e culturais brasileiras.
Podem ser trabalhados nos estudos de pensamentos, nos processos históricos, espaciais e socioculturais, reconhecendo suas diversas etnias e discutindo questões dos povos \textbf{indígenas}, quilombolas, \textbf{imigrantes} europeus e asiáticos no Brasil.
Uma oportunidade para direcionar reflexões sobre a ideia de imposição de determinadas culturas em detrimento de outras, que teria resultado na hegemonização da globalização, e evidenciar que existe um outro olhar onde existem múltiplas culturas que habitam em harmonia justamente em função da possibilidade das relações globais.
O multiculturalismo deve ser desenvolvido para criar um ambiente que aceite melhor as diferenças e assim despertar problematizações como as questões de \textbf{racismo} e preconceito, além de poder avaliar e entender o propósito cultural ou político envolvido, despertando nos estudantes o sentimento de diversidade, em que aprendam a respeitar as diferenças e que se defronte com assuntos como \textbf{identidade} cultural e de gênero.
O tema Economia envolve questões acerca do Trabalho, Educação Financeira e Educação Fiscal.
É de extrema importância para refletir sobre o universo produtivo, bem como as dinâmicas do mundo do trabalho e o preparo profissional do estudante.
Além disso, estudar economia garante o benefício de se ter indivíduos que compreendam minimamente lidar com questões econômicas e desenvolver consciência em relação às decisões do governo.
Os estudantes irão dispor de maior clareza das implicações de políticas governamentais em seus próprios rendimentos, tendo a oportunidade de se programarem para um futuro financeiro mais saudável e rentável.
Cidadania e Civismo tratam de pontos que envolvem as temáticas de Vidas Familiares e Sociais, Educação para o Trânsito, Educação em Direitos Humanos, Direitos da Criança e do Adolescente, Processo de Envelhecimento respeito e valorização do Idoso.
É um tema que aborda os valores morais do indivíduo evidenciando a importância para a sua inserção em comunidade, possibilitando o usufruto dos direitos e deveres da vida cidadã.
Promove o desenvolvimento da consciência humana, o pertencimento da comunidade, o incentivo a participação política e a formação e conscientização da prática cidadã, valorizando as características e os posicionamentos das diferentes gerações, fortalecendo o diálogo entre todos.
A temática de Ciência e Tecnologia trabalha a utilização qualificada das tecnologias e dos conteúdos das mídias para promover uma aprendizagem integradora e contribuir para o importante papel que tem a escola como ambiente de inclusão digital e de utilização crítica das tecnologias da informação e comunicação.
Deve estimular no estudante um olhar questionador e crítico e desenvolver a aptidão para solucionar situações-problema, estabelecendo conexões com o conhecimento adquirido no âmbito escolar.

% matched lemmas: antropológico, atravessar, capital, fronteira, homem, identidade, imigrante, indígena, racismo, religioso, território
\end{textsample}
